\section{\emph{Clean Code} e \emph{Effective Java}: o código}

O \citetitle{livro:martin2008} trata do código em si: nomes (cap.~2),
funções (cap.~3), comentários (cap.~4), tratamento de erros (cap.~7),
testes (cap.~9), classes (cap.~10), concorrência (cap.~13) e um catálogo
de \termo{code smells}{\enquote{Cheiros} de código: sinais de que algo
está mal estruturado, mesmo que funcione.} (cap.~17). O
\citetitle{livro:bloch2017} complementa com conselhos específicos de
Java, organizados em itens numerados. Vários deles ganham explicação
concreta no bytecode desta parte:

\small
\begin{longtable}{@{}p{5.6cm} p{6.6cm}@{}}
\toprule
Item do \emph{Effective Java} (3ª ed.) & Onde o bytecode mostra \\
\midrule \endhead
6. Avoid creating unnecessary objects & Autoboxing escondido,
  seção~\ref{sec:autoboxing} \\
9. Prefer try-with-resources to try-finally & \texttt{addSuppressed}
  gerado pelo \texttt{javac}, seção~\ref{sec:twr} \\
42. Prefer lambdas to anonymous classes & \texttt{invokedynamic} em vez
  de uma classe a mais no JAR, seção~\ref{sec:invokedynamic} \\
61. Prefer primitive types to boxed primitives & \texttt{==} entre
  \texttt{Integer}, seção~\ref{sec:autoboxing} \\
63. Beware the performance of string concatenation & Concatenação via
  \texttt{invokedynamic}, seção~\ref{sec:invokedynamic} \\
67. Optimize judiciously & JIT, warm-up e medição,
  capítulo~\ref{cap:jvm-por-dentro} \\
78. Synchronize access to shared mutable data & \texttt{contador++} são
  três instruções, seção~\ref{sec:modelo-memoria} \\
\bottomrule
\end{longtable}
\normalsize

\paragraph{Código limpo e desempenho não brigam (quase nunca).}
A JVM otimiza o código que roda muito, no momento em que ele roda
(capítulo~\ref{cap:jvm-por-dentro}). Métodos pequenos, que o
\emph{Clean Code} recomenda, são candidatos naturais a
\termo{inlining}{Otimização do compilador JIT (que compila o código
durante a execução; ver seção~\ref{sec:jit}) que copia o corpo do método
chamado para dentro do chamador, eliminando a chamada.}. Antes de sacrificar
legibilidade por desempenho, meça; o item 67 do \emph{Effective Java}
diz o mesmo.
